\textbf{结论名称：}一般式中垂直判定（\(A_1A_2+B_1B_2=0\)）

\textbf{适用条件：}两直线均用一般式表示，且各自的 \(A,B\) 不同时为零（即 \(A_1,B_1\) 不全为零，\(A_2,B_2\) 不全为零）。

\textbf{变量说明：}设 \(l_1:A_1x+B_1y+C_1=0\)，\(l_2:A_2x+B_2y+C_2=0\)，其中系数 \(A_i,B_i,C_i\in\mathbb{R}\)。

\textbf{核心公式：}
\[
\boxed{A_1A_2+B_1B_2 = 0 \iff l_1 \perp l_2}
\]

\textbf{使用场景：}无需化为斜截式即可直接判定垂直，特别适用于含有与坐标轴平行的直线（斜率不存在）的情形；常用于求垂线方程、三角形的高、对称问题等。

\textbf{简单例子：}
判断 \(l_1:2x+3y-5=0\) 与 \(l_2:3x-2y+1=0\) 是否垂直。
计算：\(A_1A_2+B_1B_2=2\times3+3\times(-2)=0\)，故 \(l_1\perp l_2\)。
又例如 \(l_3:x=2\;(1\cdot x+0\cdot y-2=0)\) 与 \(l_4:y=3\;(0\cdot x+1\cdot y-3=0)\)，有 \(1\times0+0\times1=0\)，亦垂直。